% =============================================================================
%  3-Conclusiones.tex — Conclusiones del registro de avances
% =============================================================================

\section{Conclusiones}
\label{sec:conclusiones}

% ── INSTRUCCIONES ─────────────────────────────────────────────────────────────
% Resume los hallazgos del periodo. Algunas preguntas guía:
%   • ¿Qué avances fueron más significativos?
%   • ¿Qué redefiniciones cambiaron el rumbo del proyecto?
%   • ¿Qué apartados requieren atención prioritaria?
%   • ¿Qué herramientas o decisiones resultaron acertadas?
% ─────────────────────────────────────────────────────────────────────────────

Durante el periodo comprendido entre las versiones \texttt{\VerRef}
y \texttt{\VerActual} del repositorio \texttt{\NombreRepo}, se trabajo principalmente 
en los siguientes puntos:
\begin{itemize}
    \item \textbf{Creación de Capítulo Anexo:} Se avanzó en la creación de un apartado especializado 
    para el código, descripción técnica como \textbf{LOGS}, \textbf{Fragmentos de Código} y aparados 
    que no corresponden a la estructura y eje principal (temas relacionados con el urbanismo).
    \item \textbf{Mejora en Visualización:} Se crea el capítulo llamado \textbf{Transport-gis: El motor de Visualización} 
    el cual funge como la última capa para presentar mapas y estadísticas de valores. En este contexto se usará Qgis como principal 
    y se propone usar Tableau para los datos. En cuanto a imágenes, se optó por cambiarlos al formato de \termino{TiKz} para 
    la creación de diagramas y gráficos vectoriales.
    \item \textbf{Síntesis en la Redacción:} Debido a los dos puntos anteriores, se necesitó realizar una síntesis de los contenidos 
    para mejorar la coherencia y fluidez del documento.
\end{itemize}
